\documentclass{article}

%%% Packages
\usepackage{fp}
\usepackage{booktabs}
\usepackage{numprint}
\usepackage{ifthen}
\usepackage{array}
\usepackage{xcolor}
\usepackage{threeparttable}
\usepackage{makecell}
\npthousandsep{,}

%%% Setup
\newcolumntype{R}[1]{>{\raggedleft\arraybackslash}p{#1}}
\newcolumntype{L}[1]{>{\raggedright\arraybackslash}p{#1}}

\newlength\BARSIZE  \setlength\BARSIZE{0.9cm}
\newcommand*\ChartWithAbsoluteAndPercentageNumbers[2]{
    \numprint{#1}&\ShowPercentageNoHundred{#1}{#2}&\inlinechart{#1}{#2}}

\newcommand{\inlinechart}[2]{%
\FPeval{\BLACKBARSIZE}{#1/#2}\textcolor{black!80}{\rule{\BLACKBARSIZE\BARSIZE}{1.6ex}}%
\FPeval{\BLACKBARSIZE}{1 - (#1/#2)}\textcolor{black!10}{\rule{\BLACKBARSIZE\BARSIZE}{1.6ex}}%
}

\newcommand{\ShowPercentageNoHundred}[2]{%
    \FPeval\percentage{round(#1/#2*100,1)}%
    \FPeval\percentageRounded{round(#1/#2*100,0)}%
    \ifthenelse
        {\equal{\percentageRounded}{0}}%if
        {(\hspace*{4.3pt}\FPprint{percentage}\%)}%then
        {%else
        \ifthenelse
            {\equal{\percentageRounded}{1}}%if
            {(\hspace*{4.3pt}\FPprint{percentage}\%)}%then
            {%else
            \ifthenelse
                {\equal{\percentageRounded}{2}}%if
                {(\hspace*{4.3pt}\FPprint{percentage}\%)}%then
                {%else
                \ifthenelse
                    {\equal{\percentageRounded}{3}}%if
                    {(\hspace*{4.3pt}\FPprint{percentage}\%)}%then
                    {%else
                    \ifthenelse
                        {\equal{\percentageRounded}{4}}%if
                        {(\hspace*{4.3pt}\FPprint{percentage}\%)}%then
                        {%else
                        \ifthenelse
                            {\equal{\percentageRounded}{5}}%if
                            {(\hspace*{4.3pt}\FPprint{percentage}\%)}%then
                            {%else
                            \ifthenelse
                                {\equal{\percentageRounded}{6}}%if
                                {(\hspace*{4.3pt}\FPprint{percentage}\%)}%then
                                {%else
                                \ifthenelse
                                    {\equal{\percentageRounded}{7}}%if
                                    {(\hspace*{4.3pt}\FPprint{percentage}\%)}%then
                                    {%else
                                    \ifthenelse
                                        {\equal{\percentageRounded}{8}}%if
                                        {(\hspace*{4.3pt}\FPprint{percentage}\%)}%then
                                        {%else
                                        \ifthenelse
                                            {\equal{\percentageRounded}{9}}%if
                                            {(\hspace*{4.3pt}\FPprint{percentage}\%)}%then
                                            {%else
                                            \ifthenelse
                                                {\equal{\percentage}{9.5}}%if
                                                {(\hspace*{4.3pt}\FPprint{percentage}\%)}%then
                                                {%else
                                                \ifthenelse
                                                    {\equal{\percentage}{9.6}}%if
                                                    {(\hspace*{4.3pt}\FPprint{percentage}\%)}%then
                                                    {%else
                                                    \ifthenelse
                                                        {\equal{\percentage}{9.7}}%if
                                                        {(\hspace*{4.3pt}\FPprint{percentage}\%)}%then
                                                        {%else
                                                        \ifthenelse
                                                            {\equal{\percentage}{9.8}}%if
                                                            {(\hspace*{4.3pt}\FPprint{percentage}\%)}%then
                                                            {%else
                                                            \ifthenelse
                                                                {\equal{\percentage}{9.9}}%if
                                                                {(\hspace*{4.3pt}\FPprint{percentage}\%)}%then
                                                                {%else
                                                                    (\FPprint{percentage}\%)
                                                                }%
                                                            }%
                                                        }%
                                                    }%
                                                }%
                                            }%
                                        }%
                                    }%
                                }%
                            }%
                        }%
                    }%
                }%
            }%
        }%
}%


\begin{document}

%%%%%%%%%%%%%%%%%%%%%%%%
%% Table 6
%%%%%%%%%%%%%%%%%%%%%%%%

%%% Variables
\FPeval\varTotalOfPapersWithLinkICSE{1001}
\FPeval\varTotalOfPapersWithLinkICSEFSE{1696}
\FPeval\varNumLinksNotAvailableICSEFSE{292}
\FPeval\varNumLinksPublicAvailableICSEFSE{1862}
\FPeval\varNumLinksRestrictAccessICSEFSE{17}
\FPeval\varNumPapersLinkNotAvailableICSEFSE{263}
\FPeval\varNumPapersLinkPublicAvailableICSEFSE{1469}
\FPeval\varNumPapersLinkRestrictAccessICSEFSE{16}
\FPeval\varPapersHaveOnlyNotAvailableLinksICSEFSE{217}
\FPeval\varPapersHaveOnlyPubliclyAvailableLinksICSEFSE{1419}
\FPeval\varPapersHaveOnlyRestrictAccessLinksICSEFSE{9}
\FPeval\varTotalPapersUniqueInCategoryICSEFSE{1645}
\FPeval\varTotalOfLinksICSEFSE{2171}

\setcounter{table}{5}%Table 6


\begin{table}[htbp]
\begin{threeparttable}
    \caption{Statuses of the links of ICSE and FSE papers.}
    \label{tab:status-of-the-links}
    \centering
    \small
    \begin{tabular}{@{}p{.23\textwidth} R{.06\textwidth} @{\hskip 0.02in} R{.04\textwidth} @{\hskip 0.26in} L{.08\textwidth} R{.06\textwidth} @{\hskip 0.02in} R{.04\textwidth} @{\hskip 0.26in} L{.08\textwidth} R{.06\textwidth} @{\hskip 0.02in} R{.04\textwidth} @{\hskip 0.26in} L{.08\textwidth}@{}}
    
        \toprule
    
         Status of the link & 
         \multicolumn{3}{c}{\# Links (\numprint{\varTotalOfLinksICSEFSE})$^*$} &
         \multicolumn{3}{c}{\# Papers (\numprint{\varTotalOfPapersWithLinkICSEFSE})} &
         \multicolumn{3}{c}{\thead{\# Papers unique \\ in category (\numprint{\varTotalPapersUniqueInCategoryICSEFSE})}} \\
        
        \midrule

        Publicly available &
        \ChartWithAbsoluteAndPercentageNumbers{\varNumLinksPublicAvailableICSEFSE}{\varTotalOfLinksICSEFSE} &
        \ChartWithAbsoluteAndPercentageNumbers{\varNumPapersLinkPublicAvailableICSEFSE}{\varTotalOfPapersWithLinkICSEFSE} &
        \ChartWithAbsoluteAndPercentageNumbers{\varPapersHaveOnlyPubliclyAvailableLinksICSEFSE}{\varTotalPapersUniqueInCategoryICSEFSE} \\

         Not available &
         \ChartWithAbsoluteAndPercentageNumbers{\varNumLinksNotAvailableICSEFSE}{\varTotalOfLinksICSEFSE} &
         \ChartWithAbsoluteAndPercentageNumbers{\varNumPapersLinkNotAvailableICSEFSE}{\varTotalOfPapersWithLinkICSEFSE} &
         \ChartWithAbsoluteAndPercentageNumbers{\varPapersHaveOnlyNotAvailableLinksICSEFSE}{\varTotalPapersUniqueInCategoryICSEFSE} \\

         Restricted access &
         \ChartWithAbsoluteAndPercentageNumbers{\varNumLinksRestrictAccessICSEFSE}{\varTotalOfLinksICSEFSE} &
         \ChartWithAbsoluteAndPercentageNumbers{\varNumPapersLinkRestrictAccessICSEFSE}{\varTotalOfPapersWithLinkICSEFSE} &
         \ChartWithAbsoluteAndPercentageNumbers{\varPapersHaveOnlyRestrictAccessLinksICSEFSE}{\varTotalPapersUniqueInCategoryICSEFSE} \\

        \bottomrule
        
    \end{tabular}
    \begin{tablenotes}[flushleft]
      \item [*] A paper might be counted in more than one category if it contains more than one link.
    \end{tablenotes}
  \end{threeparttable}
\end{table}

\end{document}